\chapter{Complexity — Layout2}
%%% generated by the blueprint pipeline from TCSlib.Complexity.SpaceComplexity.Codec.Layout2
\section{Overview}

Declarations of the Lean module \texttt{TCSlib.Complexity.SpaceComplexity.Codec.Layout2}.

\section{Declarations}

\begin{lemma}[Bit retrieval from nested pair encodings]
\lean{Complexity.Codec.getD_w2}
\label{Complexity.Codec.getD_w2}
\leanok
\uses{Complexity.Codec.pe_getD_fst, Complexity.Codec.pe_getD_snd, Turing.pairEncode}
\difficulty{3}
Let $x$, $w$, $w'$, and $r$ be finite bit strings, and let $i$ be a natural number with
$i<|w'|$. In the encoding obtained by pairing $x$ with the encoding of $w$ paired with
the encoding of $w'$ and $r$, the bit at position $2|x|+2|w|+4+2i$ equals the bit at
position $i$ of $w'$, with positions counted from zero.
\end{lemma}

\begin{definition}[Second register layout]
\lean{Complexity.Codec.lay2R}
\label{Complexity.Codec.lay2R}
\leanok
\uses{Complexity.Codec.Regs, Complexity.Codec.W}
Let $N$ be a finite nondeterministic Turing machine, let $R$ be a register record, and
let $b,s,wp,pos$ be natural numbers. The second register layout is obtained from $R$ by
setting $E2$ and $Q1$ to $b+4W(N)s+2W(N)$, $F2$ to that quantity plus $2wp$, $U$ to that
quantity plus $2wp+4$, $Pos2$ to $pos$, and $T$ to $0$. Every other register retains its
value from $R$.
\end{definition}

\begin{theorem}[Execution of the second layout stage]
\lean{Complexity.Codec.layout2_run}
\label{Complexity.Codec.layout2_run}
\leanok
\uses{Complexity.Codec.CA, Complexity.Codec.Reach, Complexity.Codec.Reach.cast, Complexity.Codec.Reach.single', Complexity.Codec.Reach.trans, Complexity.Codec.Regs, Complexity.Codec.Regs.Bd, Complexity.Codec.Regs.Bd.val_le, Complexity.Codec.Regs.val, Complexity.Codec.Regs.val_E2, Complexity.Codec.W, Complexity.Codec.bitsVal_lt_two_pow, Complexity.Codec.bnd, Complexity.Codec.bnd_ge, Complexity.Codec.cpDst, Complexity.Codec.cpSrc, Complexity.Codec.cp_run, Complexity.Codec.dcAcc, Complexity.Codec.dcBase, Complexity.Codec.dcPtr, Complexity.Codec.dcStride, Complexity.Codec.dc_run, Complexity.Codec.gd_of, Complexity.Codec.getD_w2, Complexity.Codec.halt_rej, Complexity.Codec.icCount, Complexity.Codec.icL, Complexity.Codec.icReg, Complexity.Codec.ic_run, Complexity.Codec.lay2R, Complexity.Codec.orc, Complexity.Codec.orc_bit, Complexity.Codec.scDst, Complexity.Codec.scFactor, Complexity.Codec.scSrc, Complexity.Codec.sc_run, Complexity.Codec.two_pow_bits_le, Complexity.Codec.wP, Complexity.Codec.wP_le, Complexity.LogProg.AHalt, Complexity.LogProg.GuardLe, Complexity.LogProg.PreS, Complexity.LogProg.astep, Complexity.LogProg.bitsVal, Complexity.LogProg.vword_whole₁, Turing.length_pairEncode, Turing.pairEncode}
\difficulty{7}
Let $N$ be a finite nondeterministic Turing machine with Boolean input symbols, and let
$y,x,w,w',r$ be Boolean words with $y$ equal to the nested pair encoding of $x$, the
pair encoding of $w$, and the pair encoding of $w'$ with remainder $r$. Let $R$ be a
register record and $s$ a natural number. Suppose every register of $R$ is at most
$(8W(N)+8)(|y|+1)$, and that $R.B=2|x|+2|w|+4$, $R.P=2|x|+2|w|+2|w'|+6$, $R.S=s\le
2|y|$, $R.Wp$ is the length of the binary representation of $|x|$, and
$R.E2=R.F2=R.T=R.Pos2=0$. Then either $|w'|=W(N)(2s+1)+R.Wp+1$ and the bit of $w'$ at
position $W(N)(2s+1)+R.Wp$ is true, in which case the step-mode program reaches the
input-checking stage from the copying stage with registers set as follows:
$E2=b+4W(N)s+2W(N)$, $F2=b+4W(N)s+2W(N)+2R.Wp$, $U=b+4W(N)s+2W(N)+2R.Wp+4$,
$Q1=b+4W(N)s+2W(N)$, $Pos2$ equal to the value of the first $R.Wp$ bits of $w'$ after
its first $W(N)(2s+1)$ bits are dropped, and $T=0$, where $b=2|x|+2|w|+4$; or this
length-and-bit condition fails, and the step-mode program halts with answer false. In
both cases, the run satisfies the register bound and safety condition at every time
before its endpoint.
\end{theorem}

\begin{definition}[Option-to-natural-number code]
\lean{Complexity.Codec.zcode}
\label{Complexity.Codec.zcode}
\leanok
Define a code for an optional Boolean by assigning $0$ to no value, $1$ to false, and
$2$ to true.
\end{definition}

\begin{theorem}[Input Symbol Simulation Run]
\lean{Complexity.Codec.insym_run}
\label{Complexity.Codec.insym_run}
\leanok
\uses{Complexity.Codec.CA, Complexity.Codec.Reach, Complexity.Codec.Reach.cast, Complexity.Codec.Reach.single', Complexity.Codec.Reach.trans, Complexity.Codec.Regs, Complexity.Codec.Regs.Bd, Complexity.Codec.Regs.Bd.val_le, Complexity.Codec.Regs.val, Complexity.Codec.bnd, Complexity.Codec.cpDst, Complexity.Codec.cpSrc, Complexity.Codec.cp_run, Complexity.Codec.gd_of, Complexity.Codec.inpOf, Complexity.Codec.length_le_bnd, Complexity.Codec.orc_bit, Complexity.Codec.pe_getD_fst, Complexity.Codec.scDst, Complexity.Codec.scFactor, Complexity.Codec.scSrc, Complexity.Codec.sc_run, Complexity.Codec.zcode, Complexity.LogProg.PreS, Complexity.LogProg.astep, Complexity.LogProg.vword_whole₁, Turing.length_pairEncode, Turing.pairEncode}
\difficulty{6}
Let $N$ be a finite nondeterministic Turing machine with Boolean input symbols, and let
$x$, $y$, and $r$ be finite Boolean words such that $y$ is the pair encoding of $x$ and
$r$. For a register record $R$, suppose every register is at most $B=(8W(N)+8)(|y|+1)$,
its position register is at most $|x|+1$, its $N_1$ register equals $|x|+1$, and its
$T$, $Z$, $Q_3$, and $Q_4$ registers are zero. Then some $u\le B$ exists such that the
step-mode program, starting at label $\mathrm{inZ}$ with register values given by $R$,
reaches label $\mathrm{clrQ1v}$ with values obtained by replacing $Z$ by the code of the
input symbol at position $R.\mathrm{Pos}$ in $x$, $U$ by $u$, and $Q_3,Q_4$ by
$R.A,R.B$; throughout the run, the program's safety condition holds and every register
remains at most $B$.
\end{theorem}
